\documentclass{article}


\usepackage{PRIMEarxiv}

\usepackage[utf8]{inputenc} % allow utf-8 input
\usepackage[T1]{fontenc}    % use 8-bit T1 fonts
\usepackage{hyperref}       % hyperlinks
\usepackage{url}            % simple URL typesetting
\usepackage{booktabs}       % professional-quality tables
\usepackage{amsfonts}       % blackboard math symbols
\usepackage{amsmath, amssymb}
\usepackage{nicefrac}       % compact symbols for 1/2, etc.
\usepackage{microtype}      % microtypography
\usepackage{graphicx}
\graphicspath{{media/}}     % organize your images and other figures under media/ folder
\usepackage{libertine}
\usepackage{libertinust1math}
\usepackage{dsfont}  % Double stroke fonts (provides \mathds)
\usepackage{nicefrac}
\usepackage{todonotes}



% Theorem environments
\usepackage{amsthm}  % For theorem-like environments
\theoremstyle{plain}
\newtheorem{theorem}{Theorem}
\newtheorem{lemma}[theorem]{Lemma}
\newtheorem{corollary}[theorem]{Corollary}
\newtheorem{proposition}[theorem]{Proposition}
\newtheorem{algorithm}[theorem]{Algorithm}
\newtheorem{assumption}[theorem]{Assumption}

\theoremstyle{definition}
\newtheorem{definition}{Definition}
\newtheorem{example}{Example}

\theoremstyle{remark}
\newtheorem{remark}{Remark}
    
\NewDocumentCommand{\br}{m}{\left( #1 \right)}
\NewDocumentCommand{\bs}{m}{\left[ #1 \right]}
\NewDocumentCommand{\bc}{m}{\left\{ #1 \right\}}

% Sets and number systems using dsfont (mathds)
\NewDocumentCommand{\R}{}{\mathds{R}}  % Real numbers
\RenewDocumentCommand{\H}{}{\mathds{H}}  % Hilbert space
\NewDocumentCommand{\N}{}{\mathds{N}}  % Natural numbers
\NewDocumentCommand{\Z}{}{\mathds{Z}}  % Integers
\NewDocumentCommand{\Q}{}{\mathds{Q}}  % Rational numbers
\NewDocumentCommand{\C}{}{\mathds{C}}  % Complex numbers
\NewDocumentCommand{\ind}{}{\mathds{1}}  % Complex numbers

% Mathematical operators and functions
\NewDocumentCommand{\defeq}{}{\coloneq} % Definition equal (
\NewDocumentCommand{\der}{}{\mathrm{d}}  % Differential oper
\NewDocumentCommand{\cder}{}{\partial}  % Differential oper
\NewDocumentCommand{\mul}{mm}{\langle #1, #2 \rangle}  % Inn
\NewDocumentCommand{\norm}{m}{\left\lVert #1 \right\rVert}  
\NewDocumentCommand{\abs}{m}{\left\lvert #1 \right\rvert}  %


\NewDocumentCommand{\maximize}{}{\text{maximize }}
\NewDocumentCommand{\minimize}{}{\text{minimize }}
\NewDocumentCommand{\subjectto}{}{\text{subject to }}
\NewDocumentCommand{\overall}{}{\text{over all }}


% Additional common mathematical abbreviations
\DeclareMathOperator*{\argmin}{arg\,min}
\DeclareMathOperator*{\argmax}{arg\,max}
\DeclareMathOperator{\dom}{dom}
\DeclareMathOperator{\epi}{epi}
\DeclareMathOperator{\prox}{prox}
\DeclareMathOperator{\conv}{conv}
\DeclareMathOperator{\tr}{tr} % Already using \operatorname,
\DeclareMathOperator{\rank}{rank} % Already using \operatorn

% More definitions as needed
\NewDocumentCommand{\grad}{}{\nabla}

\let\olddisplayopen\[
\let\olddisplayclose\]

% Redefine \[ and \]
\RenewDocumentCommand{\[}{}{%
  \begin{equation}\begin{aligned}%
}
\RenewDocumentCommand{\]}{}{%
  \end{aligned}\end{equation}%
}



%% Title
\title{Real-Time Semantic Acoustic Imaging via Variational Posterior Tracking and Generative Mixture Fitting
%%%% Cite as
%%%% Update your official citation here when published 
\thanks{\textit{\underline{Citation}}: 
\textbf{Authors. Title. Pages.... DOI:000000/11111.}} 
}

\author{
  Bogdan Jastrzębski \\
  Samsung Research Institute Poland (SRPOL) \\
  Warsaw \\
  \texttt{b.jastrzebsk@samsung.com} \\
  %% \AND
  %% Coauthor \\
  %% Affiliation \\
  %% Address \\
  %% \texttt{email} \\
  %% \And
  %% Coauthor \\
  %% Affiliation \\
  %% Address \\
  %% \texttt{email} \\
  %% \And
  %% Coauthor \\
  %% Affiliation \\
  %% Address \\
  %% \texttt{email} \\
}


\begin{document}
\maketitle


\begin{abstract}
Current semantic acoustic imaging baselines reconstruct dense $180^\circ \times 360^\circ$ spatial maps independently per frame. These systems rely on underdetermined numerical solvers that introduce checkerboard artifacts, forcing the use of early stopping which precludes temporal warm-starting~\cite{baseline2024}. We propose abandoning dense map regression in favor of a generative Bayesian framework via Variational Inference. We model the unknown true sound sources as a set of points, and the observed audio-visual spatial data as generated by a Gaussian Mixture Model (GMM) conditioned on these points. Instead of tracking rigid point estimates across time, we fit a variational approximate posterior for the current frame, and propagate this full probability distribution to the next frame to serve as its prior. This retains spatial uncertainty across time, achieving fluid real-time temporal coherence while analytically circumventing the artifacts of dense numerical solvers~\cite{baseline2024}.
\end{abstract}


% keywords can be removed
\keywords{Sound Event Localization and Detection \and Acoustic Imaging \and Variational Inference \and Gaussian Mixture Model \and Semantic Audio}


\section{Introduction}
We begin our analysis with a compact 4-microphone array, but the reasoning generalizes to an arbitrary $M$ microphones. Because baseline pipelines map $M$ microphones onto $N$ spatial directions, the resulting linear operator is fundamentally singular~\cite{baseline2024}. Exact solvers like ADMM converge to minimizers filled with ringing artifacts, requiring early stopping to produce smooth approximations~\cite{baseline2024}. Consequently, current models cannot use the solution from time $t-1$ to initialize time $t$, because warm-starting accelerates convergence toward the artifact-ridden exact minimizer~\cite{baseline2024}. We bypass this bottleneck by modeling the scene as a sparse, unknown set of point sources and fitting an approximate posterior distribution over their locations.

\section{Probabilistic Generative Model}
Let $\mathbf{P}^{(t)} = \{p_1^{(t)}, \dots, p_K^{(t)}\}$ represent the unknown spatial locations of $K$ sound sources at time $t$. The observed spatial features extracted from the audio and visual front-ends are represented as data points $\mathbf{D}^{(t)} = \{x_1, \dots, x_N\}$.

\subsection{The Likelihood (Observation Model)}
We assume the sound diffuses spatially from the true points $\mathbf{P}^{(t)}$, generating observations according to a Gaussian Mixture Model:
\begin{equation}
    p(\mathbf{D}^{(t)} \mid \mathbf{P}^{(t)}) = \prod_{i=1}^N \left( \sum_{k=1}^K \pi_k \mathcal{N}(x_i \mid p_k^{(t)}, \Sigma_A) + \pi_0 \mathcal{U}(x_i) \right)
\end{equation}
where $\Sigma_A$ represents the inherent spatial spread of the acoustic energy.

\subsection{The Temporal Prior}
The true point locations transition smoothly over time. The prior for the current frame $t$ is formed by marginalizing over the transition dynamics from the previous frame's posterior:
\begin{equation}
    p(\mathbf{P}^{(t)} \mid \mathbf{D}_{1:t-1}) = \int p(\mathbf{P}^{(t)} \mid \mathbf{P}^{(t-1)}) p(\mathbf{P}^{(t-1)} \mid \mathbf{D}_{1:t-1}) \, d\mathbf{P}^{(t-1)}
\end{equation}

\section{Variational Inference and Posterior Tracking}
Computing the exact posterior $p(\mathbf{P}^{(t)} \mid \mathbf{D}_{1:t})$ is intractable. Instead, we introduce a tractable variational distribution $q^{(t)}(\mathbf{P})$ to approximate the true posterior. We parameterize $q^{(t)}$ as a set of independent Gaussians over the point locations:
\begin{equation}
    q^{(t)}(\mathbf{P}) = \prod_{k=1}^K \mathcal{N}\left(p_k \mid \mu_k^{(t)}, \Sigma_k^{(t)}\right)
\end{equation}
Here, $\mu_k^{(t)}$ is the model's belief about the location of source $k$, and $\Sigma_k^{(t)}$ encapsulates the uncertainty of that belief.

\subsection{Optimizing the ELBO}
To fit this approximate posterior, we maximize the Evidence Lower Bound (ELBO):
\begin{equation}
    \mathcal{L}(q^{(t)}) = \mathbb{E}_{q^{(t)}}\left[ \log p(\mathbf{D}^{(t)} \mid \mathbf{P}^{(t)}) \right] - \text{KL}\left( q^{(t)}(\mathbf{P}) \parallel p(\mathbf{P}^{(t)} \mid \mathbf{D}_{1:t-1}) \right)
\end{equation}
The first term forces the approximate posterior to cover regions where the observed audio-visual data is dense (maximizing coverage). The KL-divergence term acts as a temporal regularizer, penalizing the distribution if it deviates too far from the prior established by the previous frame.

\subsection{State Propagation and Prediction}
Once the ELBO is maximized for frame $t$, we obtain our fitted approximate posterior $q^{(t)}(\mathbf{P})$. 
For the final semantic map prediction required by the evaluator, we take the MAP estimate of this distribution:
\begin{equation}
    \hat{\mathbf{P}}^{(t)} = \arg\max_{\mathbf{P}} q^{(t)}(\mathbf{P}) = \{\mu_1^{(t)}, \dots, \mu_K^{(t)}\}
\end{equation}
Crucially, the entire distribution $q^{(t)}(\mathbf{P})$ (including the variance $\Sigma_k^{(t)}$) is stored in memory and propagated to frame $t+1$ to serve as the new prior, preserving the temporal accumulation of uncertainty and enabling highly efficient, artifact-free tracking.

\section{Conclusion}
By formulating SELD as a continuous Bayesian state-estimation problem with full variational posterior tracking, the model propagates uncertainty across time instead of discarding it. This analytically circumvents checkerboard artifacts of exact numerical solvers~\cite{baseline2024} and enables fluid real-time temporal coherence by warm-starting each frame with the previous posterior distribution.

\section*{Acknowledgments}

%Bibliography
\bibliographystyle{unsrt}  
\bibliography{references}  


\end{document}
